% Auto-generated by export_appendix_responses.py --- do not edit by hand
\chapter{Respostas geradas (baseline vs condição com sinal)}
\label{AppendixB}

As respostas reproduzidas neste apêndice são outputs experimentais dos
modelos avaliados e não constituem texto autoral da dissertação.
Este apêndice apresenta as respostas geradas pela pipeline experimental,
organizadas por estímulo. Para cada estímulo indica-se o enunciado de
entrada e o par de emoções dominantes (positiva e negativa) obtido pelo
classificador GoEmotions, o mesmo sinal injectado na condição com sinal
(codificada como \textit{pipeline} nos dados). Seguem-se, por época, as respostas
\textit{baseline} (sem sinal) e da condição com sinal.
O ficheiro-fonte completo encontra-se em
\path{appendices/data/all_results.json}.

Total de células incluídas: 96.

\section{Estímulo \texttt{s01}}
\textit{I am so proud and happy about this promotion, but I am also scared and anxious that everyone will soon see that I do not deserve it.}

Ambivalência (GoEmotions): emoção positiva dominante \textit{alegria} (\texttt{joy}); emoção negativa dominante \textit{medo} (\texttt{fear}).

\subsection*{E1: DialoGPT-medium (\texttt{microsoft/DialoGPT-medium})}
\textbf{Baseline:}

I don't think you deserve it.

\textbf{Pipeline:}

I don't deserve it either. I'm scared and excited.

\subsection*{E2: BlenderBot (\texttt{facebook/blenderbot-400M-distill})}
\textbf{Baseline:}

Why do you think you don't deserve it? I'm sure you're doing a great job.

\textbf{Pipeline:}

I'm sure you'll do just fine. Just be confident in yourself and you'll be fine.

\subsection*{E3: Phi-3-mini-Instruct (\texttt{microsoft/Phi-3-mini-4k-instruct})}
\textbf{Baseline:}

It's completely normal to feel a mix of emotions when you receive a promotion. Your hard work and dedication have led to this achievement, and it's a testament to your abilities. Remember, it's okay to have doubts, but try to focus on your accomplishments and the positive impact you've made. You deserve this promotion, and it's important to believe in yourself.

\textbf{Pipeline:}

Congratulations on your promotion! It's natural to feel a mix of emotions during such a significant change. Remember, your hard work and dedication have led you to this moment, and it's okay to take some time to adjust to your new role.

\subsection*{E4: Qwen2.5-3B-Instruct (\texttt{Qwen/Qwen2.5-3B-Instruct})}
\textbf{Baseline:}

It's natural to feel both proud and anxious during a promotion. Remember, your hard work and dedication are what led to this opportunity, and it's okay to feel uncertain about how others perceive it.

\textbf{Pipeline:}

It's understandable to feel both excited and apprehensive about a promotion. Remember, deserving something isn't always about your current abilities but about the growth and effort you've shown.

\section{Estímulo \texttt{s02}}
\textit{I'm genuinely happy for my sister's wedding, yet I also feel lonely and left behind in my own life.}

Ambivalência (GoEmotions): emoção positiva dominante \textit{alegria} (\texttt{joy}); emoção negativa dominante \textit{tristeza} (\texttt{sadness}).

\subsection*{E1: DialoGPT-medium (\texttt{microsoft/DialoGPT-medium})}
\textbf{Baseline:}

I'm sorry for your loss.

\textbf{Pipeline:}

Joyous and sad.

\subsection*{E2: BlenderBot (\texttt{facebook/blenderbot-400M-distill})}
\textbf{Baseline:}

I'm sorry to hear that. Have you talked to her about how you feel about it?

\textbf{Pipeline:}

I'm sorry to hear that. I'm sure your sister is happy for you as well.

\subsection*{E3: Phi-3-mini-Instruct (\texttt{microsoft/Phi-3-mini-4k-instruct})}
\textbf{Baseline:}

I'm glad to hear that you're happy for your sister's wedding; it's a special occasion. It's completely normal to feel a bit of loneliness during such times, especially when life seems to be moving forward for others. Remember, it's okay to take a moment for yourself and reflect on your own journey. Perhaps this could be a time to focus on personal growth or to reach out to friends and family for support.

\textbf{Pipeline:}

It's wonderful to hear that you're celebrating your sister's wedding and are happy for her. It's natural to feel a mix of emotions during such significant events. Perhaps you could find comfort in sharing your feelings with friends or family, or even consider planning a special day for yourself to reflect and enjoy.

\subsection*{E4: Qwen2.5-3B-Instruct (\texttt{Qwen/Qwen2.5-3B-Instruct})}
\textbf{Baseline:}

It's common to feel a mix of emotions during such joyous occasions. Perhaps focusing on your own happiness and well-being can help balance those feelings.

\textbf{Pipeline:}

I understand that your sister's wedding is a joyous occasion, but it's also natural to feel a bit left out sometimes. How about you reach out to friends or family members who might be able to offer some companionship during this time?

\section{Estímulo \texttt{s03}}
\textit{I love living in this city and all the opportunities here, but the rent stress is making me resent every morning.}

Ambivalência (GoEmotions): emoção positiva dominante \textit{amor} (\texttt{love}); emoção negativa dominante \textit{tristeza} (\texttt{sadness}).

\subsection*{E1: DialoGPT-medium (\texttt{microsoft/DialoGPT-medium})}
\textbf{Baseline:}

I like your username though.

\textbf{Pipeline:}

Love, sadness, etc. I love the user's attitude.

\subsection*{E2: BlenderBot (\texttt{facebook/blenderbot-400M-distill})}
\textbf{Baseline:}

I'm sorry to hear that. Is there anything you can do to make it less stressful?

\textbf{Pipeline:}

I'm sorry to hear that. Do you have any hobbies to take your mind off of it?

\subsection*{E3: Phi-3-mini-Instruct (\texttt{microsoft/Phi-3-mini-4k-instruct})}
\textbf{Baseline:}

I'm sorry to hear that you're feeling stressed about rent. It's important to find a balance between enjoying your city and managing your finances. Have you considered looking for a roommate or exploring more affordable neighborhoods? Sometimes, small changes can make a big difference in your daily life.

\textbf{Pipeline:}

I'm sorry to hear that you're feeling stressed about the rent. It's great to hear that you love living in the city and appreciate the opportunities it offers. Have you considered looking into more affordable housing options or perhaps speaking with a financial advisor to help manage your expenses?

\subsection*{E4: Qwen2.5-3B-Instruct (\texttt{Qwen/Qwen2.5-3B-Instruct})}
\textbf{Baseline:}

It's understandable that the rent stress can affect your mood, but remember to prioritize self-care and possibly explore if there are any local resources or support groups that can help manage your financial stress.

\textbf{Pipeline:}

I understand that the love you have for your city is strong, but the rent stress can definitely make mornings feel tough. How about exploring any potential ways to reduce your expenses or find more affordable housing options?

\section{Estímulo \texttt{s04}}
\textit{I'm proud that I finished the marathon, although my injury still scares me and I keep thinking I pushed too hard.}

Ambivalência (GoEmotions): emoção positiva dominante \textit{alegria} (\texttt{joy}); emoção negativa dominante \textit{medo} (\texttt{fear}).

\subsection*{E1: DialoGPT-medium (\texttt{microsoft/DialoGPT-medium})}
\textbf{Baseline:}

I don't know what to tell you.

\textbf{Pipeline:}

Joyousness and fear. I'm not sure what's going on.

\subsection*{E2: BlenderBot (\texttt{facebook/blenderbot-400M-distill})}
\textbf{Baseline:}

How long did it take you to finish it? I'm sure you did great though.

\textbf{Pipeline:}

You should be proud of yourself. It takes a lot of courage to do something like that.

\subsection*{E3: Phi-3-mini-Instruct (\texttt{microsoft/Phi-3-mini-4k-instruct})}
\textbf{Baseline:}

It's truly commendable that you completed the marathon despite your injury. Remember, it's important to listen to your body and not push beyond your limits. Perhaps consider consulting with a healthcare professional to ensure you're on the right track for recovery. Your determination is admirable, but your health should always come first.

\textbf{Pipeline:}

It's truly commendable that you've completed the marathon; that's a significant achievement. It's natural to feel some apprehension about your injury, but remember that taking care of your health is paramount. Perhaps consider consulting with a healthcare professional to ensure you're on the right track for recovery.

\subsection*{E4: Qwen2.5-3B-Instruct (\texttt{Qwen/Qwen2.5-3B-Instruct})}
\textbf{Baseline:}

It's great to hear that you finished the marathon, even with the lingering concern about pushing too hard and your injury. Recovery is important, so take the time you need to heal properly.

\textbf{Pipeline:}

It's great to hear that you finished the marathon, even with the lingering fear of pushing too hard. How are you managing to cope with those thoughts?

\section{Estímulo \texttt{s05}}
\textit{I want to start this new relationship because I care about them, but I also feel anxious that I will get hurt again.}

Ambivalência (GoEmotions): emoção positiva dominante \textit{desejo} (\texttt{desire}); emoção negativa dominante \textit{nervosismo} (\texttt{nervousness}).

\subsection*{E1: DialoGPT-medium (\texttt{microsoft/DialoGPT-medium})}
\textbf{Baseline:}

I care about you.

\textbf{Pipeline:}

Desire, nervousness, anxiousness, and anxiousness respectively.

\subsection*{E2: BlenderBot (\texttt{facebook/blenderbot-400M-distill})}
\textbf{Baseline:}

Why do you feel that way? Do you think you'll be able to get through it?

\textbf{Pipeline:}

I'm sure you'll be fine. Relationships are hard, but it's worth it in the end.

\subsection*{E3: Phi-3-mini-Instruct (\texttt{microsoft/Phi-3-mini-4k-instruct})}
\textbf{Baseline:}

It's completely normal to feel anxious about getting hurt in a new relationship, especially if you've experienced it before. Remember that taking the risk of opening up to someone new can lead to rewarding connections. It's important to communicate your feelings and concerns with your partner, and to take things at a pace that feels comfortable for you. Building trust and setting healthy boundaries can help you navigate this new relationship with care and mindfulness.

\textbf{Pipeline:}

It's completely normal to feel a mix of emotions when considering a new relationship. It's important to communicate openly with your partner about your feelings and concerns. Taking things slow and building trust can help ease your anxiety. Remember, it's okay to prioritize your emotional well-being.

\subsection*{E4: Qwen2.5-3B-Instruct (\texttt{Qwen/Qwen2.5-3B-Instruct})}
\textbf{Baseline:}

It's natural to feel anxious after past hurt. Honesty with your feelings can help both of you navigate this relationship carefully.

\textbf{Pipeline:}

It's natural to feel both drawn to someone and worried about potential hurt in a new relationship. How can I help you navigate these feelings?

\section{Estímulo \texttt{s06}}
\textit{I'm grateful my parents support me financially, yet I feel ashamed and angry that I still depend on them at my age.}

Ambivalência (GoEmotions): emoção positiva dominante \textit{gratidão} (\texttt{gratitude}); emoção negativa dominante \textit{embaraço} (\texttt{embarrassment}).

\subsection*{E1: DialoGPT-medium (\texttt{microsoft/DialoGPT-medium})}
\textbf{Baseline:}

How do you feel about your parents'financial support?

\textbf{Pipeline:}

Thank you for your feedback. Thanks for your input. Thank you.

\subsection*{E2: BlenderBot (\texttt{facebook/blenderbot-400M-distill})}
\textbf{Baseline:}

I know what you mean. I'm glad my parents are supportive of me as well.

\textbf{Pipeline:}

I know what you mean. I feel the same way. I'm glad my parents are supportive.

\subsection*{E3: Phi-3-mini-Instruct (\texttt{microsoft/Phi-3-mini-4k-instruct})}
\textbf{Baseline:}

It's completely normal to feel a mix of emotions when it comes to financial support from parents. Remember, seeking help is a sign of strength, not weakness. It's important to communicate openly with your parents about your feelings and work together to find a solution that respects your independence while acknowledging their support.

\textbf{Pipeline:}

It's completely normal to feel a mix of emotions when it comes to financial support from family. Remember that it's okay to rely on your parents, and it's a sign of their love and care for you. It's also important to communicate openly with them about your feelings and work together to find a path to financial independence that feels right for you.

\subsection*{E4: Qwen2.5-3B-Instruct (\texttt{Qwen/Qwen2.5-3B-Instruct})}
\textbf{Baseline:}

It's common to feel a mix of emotions about financial dependence, but it's important to recognize that your parents' support is a reflection of their love and care for you. You might consider talking to them about your feelings to find a way to move forward together.

\textbf{Pipeline:}

It's common to feel a mix of emotions when you're still relying on your parents financially. Gratitude for their support is a positive feeling, but the shame and anger you're experiencing are understandable too.

\section{Estímulo \texttt{s07}}
\textit{I feel relieved and even joyful that I finally submitted my thesis, but I am also disappointed and worried because parts of it are still weak.}

Ambivalência (GoEmotions): emoção positiva dominante \textit{alegria} (\texttt{joy}); emoção negativa dominante \textit{nervosismo} (\texttt{nervousness}).

\subsection*{E1: DialoGPT-medium (\texttt{microsoft/DialoGPT-medium})}
\textbf{Baseline:}

I don't know what I'm doing.

\textbf{Pipeline:}

Joyous and nervous. I'm not sure what's going on.

\subsection*{E2: BlenderBot (\texttt{facebook/blenderbot-400M-distill})}
\textbf{Baseline:}

I'm sure you did great! I know how you feel though, it can be nerve wracking.

\textbf{Pipeline:}

I'm sure you did great! I know how you feel though, it can be nerve wracking.

\subsection*{E3: Phi-3-mini-Instruct (\texttt{microsoft/Phi-3-mini-4k-instruct})}
\textbf{Baseline:}

I'm glad to hear that you've submitted your thesis and are feeling relieved and joyful about this accomplishment. It's completely normal to have mixed feelings after such a significant milestone. Remember, it's okay to feel disappointed and worried about the parts that need improvement. This is a great opportunity to refine your work and strengthen your arguments. Keep focusing on the positive progress you've made, and don't hesitate to seek feedback to enhance your thesis further.

\textbf{Pipeline:}

I'm really proud of you for submitting your thesis! It's a significant accomplishment, and it's normal to feel a mix of emotions. Remember, it's okay to take some time to review and strengthen the parts you're not as confident about. You've got this!

\subsection*{E4: Qwen2.5-3B-Instruct (\texttt{Qwen/Qwen2.5-3B-Instruct})}
\textbf{Baseline:}

It's great to hear you feel relieved and joyful about submitting your thesis. It's natural to feel disappointed and worried about the areas that need improvement, but remember, this is a normal part of the process.

\textbf{Pipeline:}

It's great to hear that you feel relieved and joyful about submitting your thesis. However, it's understandable to feel a bit nervous and disappointed about the parts that still need improvement.

\section{Estímulo \texttt{s08}}
\textit{I'm excited about moving abroad for work, while at the same time I feel sad about leaving my friends and family.}

Ambivalência (GoEmotions): emoção positiva dominante \textit{excitação} (\texttt{excitement}); emoção negativa dominante \textit{tristeza} (\texttt{sadness}).

\subsection*{E1: DialoGPT-medium (\texttt{microsoft/DialoGPT-medium})}
\textbf{Baseline:}

I don't know what to do with my life.

\textbf{Pipeline:}

I don't know what to do with my hands!!!

\subsection*{E2: BlenderBot (\texttt{facebook/blenderbot-400M-distill})}
\textbf{Baseline:}

I'm sure you'll make new friends when you get back. Where are you moving to?

\textbf{Pipeline:}

That's understandable. I'm sure you'll make new friends though. Where are you moving to?

\subsection*{E3: Phi-3-mini-Instruct (\texttt{microsoft/Phi-3-mini-4k-instruct})}
\textbf{Baseline:}

It's completely normal to feel a mix of emotions when moving abroad for work. The excitement of new opportunities and experiences can be very compelling, but it's also understandable to feel sad about leaving loved ones behind. Remember to stay connected with your friends and family through regular calls or video chats, and consider planning visits to keep those relationships strong.

\textbf{Pipeline:}

It's completely normal to feel a mix of emotions when facing such a significant life change. It's great that you're excited about the new opportunities that await you abroad. Remember, it's also important to cherish the time you have with your loved ones. Keep in touch with them, and perhaps plan a visit back home soon.

\subsection*{E4: Qwen2.5-3B-Instruct (\texttt{Qwen/Qwen2.5-3B-Instruct})}
\textbf{Baseline:}

It's natural to feel excited about new opportunities while also sad to leave behind loved ones. Balancing these emotions can be challenging, but it's a common experience when making big life changes.

\textbf{Pipeline:}

Moving abroad can be incredibly exciting, but it's also natural to feel sad about leaving behind the people you love. How are you planning to stay connected with your friends and family back home?

\section{Estímulo \texttt{s09}}
\textit{I admire my colleague's confidence in meetings, but I also feel annoyed and small whenever they interrupt me.}

Ambivalência (GoEmotions): emoção positiva dominante \textit{admiração} (\texttt{admiration}); emoção negativa dominante \textit{irritação} (\texttt{annoyance}).

\subsection*{E1: DialoGPT-medium (\texttt{microsoft/DialoGPT-medium})}
\textbf{Baseline:}

I'm with you, I'm not.

\textbf{Pipeline:}

I don't know what to do with my hands!!!

\subsection*{E2: BlenderBot (\texttt{facebook/blenderbot-400M-distill})}
\textbf{Baseline:}

Why do you feel that way? Is it because you are intimidated by them?

\textbf{Pipeline:}

I know what you mean. It's hard not to feel that way sometimes. Have you tried talking to them about it?

\subsection*{E3: Phi-3-mini-Instruct (\texttt{microsoft/Phi-3-mini-4k-instruct})}
\textbf{Baseline:}

I understand that it can be frustrating when someone interrupts you, especially in a professional setting. It's important to remember that everyone has their own communication style, and your colleague's confidence might simply be a part of their personality. Perhaps you could address this issue directly with them in a private conversation, expressing how their interruptions make you feel. This way, you can work together to find a more respectful way to communicate during meetings.

\textbf{Pipeline:}

It's great to hear that you appreciate your colleague's confidence; it's a valuable trait in any professional setting. However, it's important to feel respected and heard during meetings. Perhaps you could suggest a meeting etiquette to the team or speak to your colleague privately about how their interruptions make you feel.

\subsection*{E4: Qwen2.5-3B-Instruct (\texttt{Qwen/Qwen2.5-3B-Instruct})}
\textbf{Baseline:}

It's great that you admire your colleague's confidence, but it's important to address the interruptions as they can make you feel small and disrespected. You might consider politely reminding them of the importance of listening to others as well.

\textbf{Pipeline:}

It sounds like you're in a challenging situation where you respect your colleague's assertiveness but find it frustrating. How about discussing this with your colleague to find a way to balance their confidence with your need for respect and space?

\section{Estímulo \texttt{s10}}
\textit{I'm hopeful the medical results will be fine, yet every notification on my phone fills me with dread.}

Ambivalência (GoEmotions): emoção positiva dominante \textit{optimismo} (\texttt{optimism}); emoção negativa dominante \textit{medo} (\texttt{fear}).

\subsection*{E1: DialoGPT-medium (\texttt{microsoft/DialoGPT-medium})}
\textbf{Baseline:}

I'm sorry to hear that.

\textbf{Pipeline:}

I don't know what to do with my hands!!!

\subsection*{E2: BlenderBot (\texttt{facebook/blenderbot-400M-distill})}
\textbf{Baseline:}

I'm sorry to hear that. I hope everything turns out well for you. What are you worried about?

\textbf{Pipeline:}

I'm sorry to hear that. I hope everything turns out well for you. What are you worried about?

\subsection*{E3: Phi-3-mini-Instruct (\texttt{microsoft/Phi-3-mini-4k-instruct})}
\textbf{Baseline:}

I understand that waiting for medical results can be a stressful time. It's natural to feel anxious about the notifications, but remember that it's important to stay positive and hopeful. Perhaps you could try to distract yourself with activities you enjoy or talk to someone close to you for support. Remember, whatever the results, you're not alone, and there are people ready to help you through this.

\textbf{Pipeline:}

I understand that waiting for medical results can be a stressful time, and it's natural to feel a mix of hope and anxiety. Remember that regardless of the outcome, you have support systems in place to help you through whatever comes next. It's important to focus on the steps you can take to maintain your well-being during this period.

\subsection*{E4: Qwen2.5-3B-Instruct (\texttt{Qwen/Qwen2.5-3B-Instruct})}
\textbf{Baseline:}

I understand your concern, and it's completely normal to feel anxious during this time. It's important to stay positive and focus on the support systems around you.

\textbf{Pipeline:}

I understand that you're hopeful for the best but feel dread with each notification. How about we focus on staying positive and hopeful until you hear the results?

\section{Estímulo \texttt{s11}}
\textit{I enjoy remote work and the freedom it gives me, but I miss the office and feel disconnected from the team.}

Ambivalência (GoEmotions): emoção positiva dominante \textit{alegria} (\texttt{joy}); emoção negativa dominante \textit{tristeza} (\texttt{sadness}).

\subsection*{E1: DialoGPT-medium (\texttt{microsoft/DialoGPT-medium})}
\textbf{Baseline:}

I miss the office and the freedom it gives me.

\textbf{Pipeline:}

Joyous and sorrowful. I'm not sure what's going on.

\subsection*{E2: BlenderBot (\texttt{facebook/blenderbot-400M-distill})}
\textbf{Baseline:}

I'm sorry to hear that. What do you do in your spare time to keep yourself occupied?

\textbf{Pipeline:}

I'm sorry to hear that. What do you do in your spare time to keep yourself occupied?

\subsection*{E3: Phi-3-mini-Instruct (\texttt{microsoft/Phi-3-mini-4k-instruct})}
\textbf{Baseline:}

I understand that remote work offers flexibility and comfort, but it's also important to maintain a sense of connection with your team. Perhaps you could schedule regular video calls or virtual team-building activities to bridge the gap. Additionally, consider setting up a virtual office space where you can interact with colleagues as if you were in the same physical location.

\textbf{Pipeline:}

I understand that remote work offers you a lot of freedom, which is great. It's natural to miss the camaraderie and connection that comes with being in the office. Perhaps you could schedule regular virtual meet-ups or team-building activities to help bridge that gap and maintain a sense of community.

\subsection*{E4: Qwen2.5-3B-Instruct (\texttt{Qwen/Qwen2.5-3B-Instruct})}
\textbf{Baseline:}

It's common to feel that way, but remember that maintaining regular communication and team-building activities can help bridge the gap and keep you connected.

\textbf{Pipeline:}

I understand that remote work offers a lot of flexibility, but it can be tough to stay connected with your team and miss the social aspects of being in the office.

\section{Estímulo \texttt{s12}}
\textit{I'm glad I told the truth in that argument, although I regret how harsh I sounded and keep replaying it.}

Ambivalência (GoEmotions): emoção positiva dominante \textit{gratidão} (\texttt{gratitude}); emoção negativa dominante \textit{remorso} (\texttt{remorse}).

\subsection*{E1: DialoGPT-medium (\texttt{microsoft/DialoGPT-medium})}
\textbf{Baseline:}

I don't know what I'm doing.

\textbf{Pipeline:}

Thank you for your time. I'm sorry for the inconvenience.

\subsection*{E2: BlenderBot (\texttt{facebook/blenderbot-400M-distill})}
\textbf{Baseline:}

I'm sorry you had to go through that. I'm sure it wasn't your fault.

\textbf{Pipeline:}

I am glad you are able to look at the positive side of things. I am sure you learned from your mistake.

\subsection*{E3: Phi-3-mini-Instruct (\texttt{microsoft/Phi-3-mini-4k-instruct})}
\textbf{Baseline:}

It's commendable that you value honesty in your interactions. Remember, it's natural to reflect on our words, especially after a heated exchange. Perhaps consider expressing your thoughts in a calm and constructive manner next time. This way, you can maintain your integrity while also being mindful of how you communicate.

\textbf{Pipeline:}

It's commendable that you value honesty in your interactions. Reflecting on your words can be a powerful step towards personal growth. Remember, it's okay to express yourself, but also consider the impact your words may have on others.

\subsection*{E4: Qwen2.5-3B-Instruct (\texttt{Qwen/Qwen2.5-3B-Instruct})}
\textbf{Baseline:}

It's good that you told the truth, even if it made the situation more difficult. Moving forward, you might consider how to express your feelings more calmly next time.

\textbf{Pipeline:}

It's good to know you stood by your truth, even if it wasn't easy to express.

